% Generated draft: requires mathematical and source review.
\documentclass[11pt]{article}
\usepackage[margin=25mm]{geometry}
\usepackage{fontspec}
\setmainfont{Latin Modern Roman}
\usepackage{amsmath,amssymb,mathtools,unicode-math}
\setmathfont{Latin Modern Math}
\usepackage{enumitem,needspace,xurl,hyperref}
\setlength{\parindent}{0pt}
\setlength{\parskip}{5pt}
\setlength{\emergencystretch}{3em}
\newcommand{\N}{\mathbb{N}}
\newcommand{\Z}{\mathbb{Z}}
\newcommand{\Q}{\mathbb{Q}}
\newcommand{\R}{\mathbb{R}}
\newcommand{\C}{\mathbb{C}}
\newcommand{\F}{\mathbb{F}}
\newcommand{\A}{\mathbb{A}}
\newcommand{\PP}{\mathbb{P}}
\newcommand{\E}{\mathbb{E}}
\newcommand{\eps}{\varepsilon}
\newcommand{\dd}{\,\mathrm{d}}
\begin{document}
\section*{UnsolvedMath ID: 20000346}
\textbf{Source problem code:} catalogue entry
\subsection{Definitions and mathematical statement}
Farkas's AIM question asks, for large $g$, the minimum genus of a curve passing through a general point of the open moduli space $M_g$, and asks the analogous question for the moduli stack $\mathcal M_g$. The supplied verification explicitly identifies $M_g$, not $\overline M_g$, as the source notation. For an integral curve in a moduli space, its geometric genus can mean the genus of its smooth projective normalization; if only complete curves contained in the open space are allowed, that is an additional restriction. For a stack one must specify maps from ordinary curves versus twisted curves, whose stack structure permits inertia at marked points. These choices are not interchangeable and the source does not fix them.

The source says the stack answer is zero for low $g$, without a precise range or completeness convention. It also proposes the question for moduli $M_H$ of canonically polarized manifolds. Its motivation is to make $M_g$ non-uniruled by taking $g$ large; it separately mentions the absence of comparable conditions ensuring non-unirationality for $M_H$. No equivalence between those two properties is asserted.

For the distinct projective coarse-space covering-genus interpretation, the report claims minimum genus zero for $2\leq g\leq15$ and a lower bound of two for $g\geq22$. This does not give the exact minimum for large $g$, settle the intervening range, or answer all readings of the open-space/stack question.

Its additional necessary condition concerns a nonconstant map from a smooth projective normalized coarse curve $C$ to $\overline M_g$, with generic point in the automorphism-free interior. Put $S=f^{-1}(\overline M_g\setminus M_g)_{\mathrm{red}}$. Take a connected Galois full-level lift $\pi:C'\to C$ of degree $d$, and set $S'=\pi^{-1}(S)_{\mathrm{red}}$. For each interior branch point $p$, let $e_p$ be the ramification index of this normalized pullback. These are contact-adjusted inertia indices, not automatically the full ambient stabilizer orders. The report claims the strict log-orbifold inequality
\[
 2g(C)-2+|S|+\sum_{p\in C\setminus S}(1-1/e_p)>0,
\]
with only finitely many nonzero terms. Its stated mechanism is the identity
\[
 \frac{2g(C')-2+|S'|}{d}
 =2g(C)-2+|S|+\sum_{p\in C\setminus S}(1-1/e_p),
\]
so boundary ramification and deleted lifts combine to one reduced puncture for each boundary point. Positivity uses the nonisotrivial smooth family after level lifting in the intended characteristic-zero setting; the exact level and field hypotheses are not supplied.

When $S=\varnothing$, the report gives the quantitative bounds for the ramification divisor $R_\pi$:
\[
 g(C)=0\Longrightarrow\deg R_\pi\geq2d+2,
 \qquad g(C)=1\Longrightarrow\deg R_\pi\geq2.
\]
These are necessary constraints on a suitable level lift, not existence theorems or exact minimum-genus formulas. The generic automorphism-free condition also restricts which moduli loci the inequality can address.
\subsection{Short English statement}
Find the least geometric genus of curves through a general point of $M_g$, distinguishing open versus compactified spaces, complete versus punctured curves, and ordinary versus twisted stack maps. The report supplies projective covering-genus bounds and a necessary level-cover ramification inequality, but does not determine the original large-genus minimum in all these settings.
\subsection{Sources}
\begin{itemize}
\item Local registry: \url{data/unsolvedmath/problems.json}, permanent ID $20000346$, source code \texttt{catalogue entry}; original question, supplied open-space verification and complete research summary.
\item Michael A. van Opstall's AIM list ``Open Problems in Compact Moduli Spaces and Birational Geometry'', section 5.2 ``Other questions'', Farkas's item: \url{https://aimath.org/WWN/birational/birational.pdf}. Canonical location \texttt{aim-algebraic-geometry-notes.json}, \texttt{notes[345]}; source item 5 has extraction-broken numbering. The page footer and following ``Moduli of abelian varieties'' heading are not part of the mathematical question. No PDF was opened here.
\item The supplied verification says the displayed source is $M_g$ rather than $\overline M_g$. No full bibliographic references are supplied for the covering-genus ranges or the level-family positivity input; original literature field empty.
\item Supplied machine-generated research attempt 1: partial result, problem open at the run and currently partially solved; rigor metadata ``proved'', novelty confidence low. The non-peer-reviewed synthesis and log-orbifold obstruction require independent verification before citation.
\end{itemize}
\end{document}
